Practical use of artificial intelligence systems requires more than training a model: the system must also execute the model on user requests under latency, throughput, and availability constraints. This stage of model operation is commonly referred to as inference. Modern distributed inference of large language models based on the Mixture-of-Experts (MoE) architecture is a complex systems problem. Many existing solutions rely on collective communication and primarily optimize throughput for a fixed set of participants~\cite{aminabadi2022deepspeed,deepseekDeepEP}. However, this scheme is poorly suited for dynamic load balancing across expert blocks: request redistribution, use of expert replicas, and adding or excluding workers at runtime require a more flexible interaction model.

This research proposes an alternative client-server architecture for interaction between the main model and expert blocks. A prototype system was developed that supports asynchronous request dispatching, expert block state monitoring, dynamic selection of available expert blocks, and runtime handling of partial failures. To reduce the overhead introduced by replacing collective operations with direct interaction with expert blocks, several optimizations aimed at improving GPU computation and data transfer efficiency were implemented.

The experiments show that the proposed approach enables dynamic load balancing across expert blocks and preserves system availability under partial failures by selecting available expert replicas. The obtained results confirm the potential of the client-server approach for building flexible distributed inference systems for MoE models~\cite{shazeer2017moe,deepseekDeepEP}.
